\properlane{the-judge}{The Judge whom the faithful long to see}

\subsection{Two ways of hearing a warning}

A person who hears that the Son of Man is coming might reasonably ask for more time. At Mass he has just acknowledged his sins in the Collect. Yet the first thing Jesus says directly to the disciples in this Gospel is that they may lift their heads. The change of addressee matters. He has spoken \emph{of} people fainting for fear. He now speaks \emph{to} his friends. What allows the second group to hear a promise of deliverance where the first hear a warning of catastrophe?

Gregory the Great takes the disciples' raised heads as a sign of what they can do with their minds. The head directs the members of the body. So too the mind can direct the thoughts that crowd upon it when much of what it counted on has fallen away. Gregory was preaching to people who had lost neighbors and seen houses and churches destroyed. He recalled a storm in which people who had made plans for the next day died in the night. His hearers did not need instruction in how to picture a disaster. What they might lack was the ability to look beyond it. As he reads the Gospel, the disciples can rejoice because they will come face to face with the one they love. The end of the present order will cause grief to someone who has planted his heart so firmly in its goods that he desires no other life.\footnote{Gregory the Great, \work{Homiliae in Evangelia}, I, Homily 1, §§1--3, 5, on Lk 21:25--33.}

A man can be heartily tired of the labor of a journey and still dread reaching its end. Gregory asks how that could be if the traveler desired his home. Weariness with things as they are gives no assurance about what one wants in their place. A person whose neighborhood has been devastated might want first of all to have back what he has lost. Gregory has something more particular in mind when he says the disciples can look up. They recognize in the coming Judge the person for whose sake they undertook the journey. The reason to rejoice is the prospect of seeing him. The Roman hearers' losses explain the urgency with which Gregory pressed this question. Luke's promise allows no inference that a particular storm is the appointed sign of the end.\footnote{Gregory, \work{Homiliae in Evangelia} I.1.3, 5. Lk 21:8, 25--28.}

Jesus tells the disciples to lift their heads. In the Introit the singer has already lifted his soul and tells God to whom he has turned. Augustine hears Christ speaking in the person of his Church. She remembers the shame that came from trusting herself. Now she can tell God that she trusts him even while her enemies are watching. The enemies she fears in Augustine's exposition are quite capable of flattering her, praising the very choices that bring her low. Their approval will give her no way out of shame. As she prays the verse about learning the Lord's paths, the Church can acknowledge how much she needs to unlearn.\footnote{Augustine, \work{Enarrationes}, on Ps 24, §§1--5, English Psalm XXV. The psalmist's petition is Ps 24(25):1--4.}

Before the Gospel the same congregation asks to be shown mercy. Augustine hears in that request for salvation a desire for Christ himself. The man who receives this gift has no grounds to claim he secured his own rescue. What gives him joy is that the Savior has turned toward him and can turn him back to himself. In the Collect the Church has already asked precisely for the coming of the One she calls upon, and admitted the origin of her peril. Hearing the Alleluia after such an admission allows someone who would be afraid to see his life exposed to ask quite sincerely to meet his Judge. He can want to face the Lord from whom he seeks forgiveness.\footnote{Augustine, \work{Enarrationes}, on Ps 84, §6, at the petition of Ps 84:8 (85:7).}

\subsection{The years since first believing}

To call the coming of Christ one's salvation may seem easy for the disciples who know that he will redeem them. It is a more searching thing to hear Paul say that salvation has drawn nearer since these Christians believed. Chrysostom notices an obstacle inside that comparison. Someone may remember how eager he was in the early days of his faith and how much less he does now. As the years have gone by, the end of his earthly life has indeed drawn closer. His love may actually have grown colder. When Paul calls for waking he speaks to someone capable of falling asleep after beginning well. There would have been little point in sounding the alarm if the mere lapse of time did the work of preparing the sleeper.\footnote{John Chrysostom, \work{Homilies on Romans}, Homily 24, on Rom 13:11.}

Chrysostom draws on the racecourse. Toward the finish a runner can see the prize more clearly and so makes a greater effort. Someone who grows more slack exactly as the finish approaches is behaving in a way even a spectator at the games can recognize as strange. What Paul names at the finish explains why he can hope to wake someone. He might have told the Romans that their judgment is nearer. To say \emph{salvation} to people who still have sins to put away is to address them as people who may yet receive the prize. The word asks them to hope enough to act. Whether that day will mean salvation for someone matters very much to Chrysostom; he attaches the promise to the readiness Paul requires.

There is also the sound in the street before it is light. Chrysostom describes the stirring that begins when people hear the early bird and call one another to get up. At that hour one might perfectly well observe that it is still dark. The sleepers who heard the bird know enough to dress for day. They do not wait for full daylight to establish beyond question that their work must begin. Paul likewise tells the Romans what kind of conduct fits the approaching day. In the Gradual the assembly asks God to teach his paths just after Paul has told them how to walk. His list of the night's deeds may be thoroughly familiar. Knowing that an injury is being kept alive by one's own willingness to quarrel takes another act of recognition, as does leaving it unavenged. Someone who asks to be taught a path can begin at the point where Paul has pointed to a real departure from it.\footnote{Chrysostom, Homily 24 on Romans, on Rom 13:12--14; Augustine, \work{Enarrationes}, on Ps 24, §4.}

The last pair in the list is especially liable to escape the Christian making that judgment. Drunkenness can leave an unmistakable trace. Contention can thrive while a man believes that he has been defending a perfectly good cause. Envy can go through life under the name of fairness. The question Chrysostom poses at Paul's last command concerns what can actually be seen of Christ in such a person. He compares putting on Christ to a familiar way of speaking about friends who are always in one another's company. One learns what matters to one friend by observing the other. Gentleness toward someone with whom one has had cause to quarrel can make the company one keeps recognizable to a third person.\footnote{Chrysostom, Homily 24 on Romans, on Rom 13:13--14.}

\subsection{The fruit one can recognize}

The Communion chant says first what God will give and then what the earth will yield. Augustine distinguishes the words of the preceding verse from those actually sung. At Ps 84:12 truth springs from earth. He can say this of Christ taking mortal flesh from Mary. He can also hear in it the sinner's confession: a person made of earth at last speaks the truth about what he has done. At the next verse Augustine addresses the people listening to his sermon. The earth that is to bear its increase is already under cultivation. Confession has opened the furrows and the people have heard the word sown in them. There is still need of God's rain. It is possible to have heard exactly what one did wrong, to know it and to say so truthfully, and still find oneself taking pleasure in it.\footnote{Augustine, \work{Enarrationes}, on Ps 84, §§10--12 (English Psalm LXXXV), on Ps 84:12--13 (85:11--12).}

The chant promises that the Lord will give ``goodness,'' \emph{benignitatem}. In the English of Augustine's sermon the promise reads, ``the Lord shall give sweetness.'' Augustine asks his hearers to consider what it feels like to do justice. A man once enjoyed stealing what belonged to someone else. Now it gives him pleasure to share what he owns, and he goes looking for someone who needs it. Augustine hears the promised sweetness in that new delight.\footnote{Augustine, \work{Enarrationes}, English Psalm LXXXV, §12. Compare the 1962 Missal, no. 9, Ps 84:13 (85:12).} A Christian singing of the Lord's goodness at Communion can recognize that gift in the pleasure he now takes in helping a neighbor. If the neighbor once had reason to hide his goods from him, he too can tell what has changed.

Where, then, does a man go who needs just such a change of heart? At the Offertory he sings that he has put his trust in God. There are gifts before him, but the prayer over them assigns the power to cleanse to the holy things, and asks for their recipients to become purer. The thanksgiving that follows requires him to give voice to the direction of his own heart. He answers the call to lift it to the Lord. The priest can then speak of the majesty of the Son as equal to the Father's and join the praising angels. In Luke the Son's majesty has occasioned fear. In the Preface it is possible to acclaim that same majesty while still awaiting his appearing. A person who needs cleansing can take part in the answer of the congregation and receive the gift for which the Secret asks.

\subsection{What will last through the change of seasons}

On Gregory's account the fig tree has a disquieting fruit. What the world produces in its ripeness is its own collapse. Its growth has never guaranteed permanence. To those who have set their hope on the kingdom, however, the signs have the other implication Jesus drew from the trees. Summer is near. Gregory pictures a light in which the clouds of grief have passed. A damaged tree or a lost house in his hearers' city had ordinarily meant less shelter. Jesus' promise enables them to hope for a home in which what they love cannot be lost. The Postcommunion puts the praying Church in a temple asking for mercy while the festivals of her redemption approach. The coming festivities give her an immediate object for careful preparation. The last coming for which Gregory has his hearers prepare exceeds any celebration they can keep in this temple.\footnote{Gregory, \work{Homiliae in Evangelia} I.1.3. 1962 Missal, no. 10.}

There remains Jesus' solemn statement about the generation that will not pass until all comes to pass. Bede offers two possibilities for its extent. It can mean the whole human race, or the Jewish people in particular. Either can continue through the lifetimes of many hearers. Bede remarks on the purpose of the Amen: Truth has chosen to speak with the force of an oath so that those half asleep will attend to what they hear. Gregory quotes the saying and takes up the next verse. Speech seems among the least durable things one can perceive. The sound vanishes as the word is heard. Yet Christ's words express a judgment that will stand when the apparently solid earth under the hearers' feet has passed away. He can trust the promise about his redemption because of who has spoken it.\footnote{Bede, \work{In Lucam} VI, on Lk 21:32, PL 92, col. 590B--C. Gregory, \work{Homiliae in Evangelia} I.1.4, on Lk 21:32--33.}

\begin{fourSenses}
\item[Literal.] The psalmist calls upon God under threat. Paul bids the Romans leave behind practices that thrive in the dark. Jesus tells his disciples what the approach of the Son of Man will mean for them while the nations are in distress.
\item[Allegorical.] In Augustine's account Christ lends his voice to his Church so that a people brought low can lift their soul to God. The Son they ask to come to their aid is the Son whose majesty they adore with the Father and the Spirit.
\item[Moral.] Someone who would meet Christ gladly can begin by giving up the pleasure of making an adversary feel an old injury. Paul's words about contention and Augustine's example of a thief who learns to enjoy giving allow the change to be tested in the good another person receives.
\item[Anagogical.] The faithful hope for the meeting with God that Gregory likens to summer after clouds of grief. The promise of that meeting rests on words which will endure when the visible heavens have passed.
\end{fourSenses}
